\subsection{扭转的复合}\label{sec:5e}
取 c-有限、三角化的 $\Ai$ 范畴 $\A$，固定对象 $Y_1,\ldots,Y_m$。从任意对象 $X=X_m$ 出发，依次进行扭转，令 $X_{k-1}=T_{Y_k}X_k$。由定义，有正合三角
\SourceDiagram{diagram-5e}{0.2990766}
\SourcePage{78}
其中 $Z_k=\Hom_{H(\A)}(Y_k,X_k)$。反复应用\stref{推论}{3.19}，得到合成的正合三角
\SourceEquationDiagram{diagram-5-12}{0.30946121}{5.12}
其中对象 $Y$ 通过一系列正合三角由 $Z_k\otimes Y_k$ 构造而成，因此属于 $\A$ 中由 $Y_k$ 生成的三角化 $\Ai$ 子范畴。现在假设 $[i_1]\circ\cdots\circ[i_m]=0$。\znote{70}{在三角范畴中，正合三角 $X\to Y\to Z\to TX$ 分裂（即 $Y\cong X\oplus Z$）当且仅当中段复合 $X\to Y\to Z$ 为零。这里 $[i_1]\circ\cdots\circ[i_m]$ 正是合成三角 (5.12) 中从 $X$ 出发的那一段，故它为零时 (5.12) 分裂，得到 (5.13)；其中 $[-1]$ 来自第三项与锥的平移约定。[来源：nLab "triangulated category"，https:/\allowbreak /\allowbreak ncatlab.org/\allowbreak nlab/\allowbreak show/\allowbreak triangulated+\allowbreak category（2026-11-15 访问）；并见 Seidel, Fukaya categories and Picard–Lefschetz theory, 笔记 3.18。]}那么\eref{5.12}按\stref{笔记}{3.18}所述方式分裂，故
\begin{equation}\label{eq:5.13}
X\oplus T_{Y_1}\cdots T_{Y_m}(X)[-1]\cong Y.
\end{equation}
显然，这个观察在具体情形中是否有用，取决于能否证明 $[i_k]$ 的复合为零。下面给出两个特别简单的情形：
\begin{corollary}\label{st:5.7}
设 $\A$ 是 c-有限、三角化的 $\Ai$ 范畴。假设存在 $Y_1,\ldots,Y_m$，使对所有对象 $X$，
\begin{equation}\label{eq:5.14}
T_{Y_1}\cdots T_{Y_m}X\cong0
\end{equation}
都是 $H^0(\A)$ 的零对象（等价地，$T_{Y_1}\cdots T_{Y_m}X$ 在该范畴中的自同态群为零）。那么 $Y_k$ 生成 $\A$，而且这个范畴分裂闭。
\end{corollary}
\begin{proof}
第一项断言直接由\eref{5.13}得到。对于第二项，考虑分裂闭包 $\Pi\A$；它显然仍 c-有限，并由\stref{引理}{4.8}三角化。由扭转的定义可见，$T_{Y_k}$ 保持直和，因此\eref{5.14}对 $\Pi\A$ 的所有对象 $X$ 都成立。于是 $Y_k$ 也生成 $\Pi\A$，这意味着嵌入 $\A\hookrightarrow\Pi\A$ 是拟等价。
\end{proof}
\begin{corollary}\label{st:5.8}
设 $\A$ 是 c-有限、三角化且分裂闭的 $\Ai$ 范畴。假设存在 $Y_1,\ldots,Y_m$ 与 $\sigma\ne0$，使得对所有 $X$，在 $H^0(\A)$ 中有
\[
T_{Y_1}\cdots T_{Y_m}(X)\cong S^{\sigma}X.
\]
那么 $Y_k$ 分裂生成 $\A$。
\end{corollary}
\begin{proof}
有限维假设意味着，对每个 $X$，存在 $s>0$，使 $\Hom^{s\sigma}_{H(\A)}(X,X)=0$。将给定对象族 $Y_k$ 与自身连接 $s$ 次。在所得正合三角\eref{5.12}中，态射
\[
X\longrightarrow(T_{Y_1}\cdots T_{Y_m})^s(X)\cong S^{s\sigma}X
\]
必然为零，因此由\eref{5.13}，$X$ 由 $Y_k$ 分裂生成。
\end{proof}

\SourcePage{79}
\subsection{扭曲复形}\label{sec:5f}
设 $\B$ 是严格含幺的 $\Ai$ 范畴，其链层次的 $\hom$ 空间有限维，并取 $\A=\Tw\B$。在此情形下，扭转函子有显式代表。具体地，给定 $Y_0,Y_1\in\Ob\A$，取张量积 $\hom_{\A}(Y_0,Y_1)\otimes Y_0$，其中如\secref{3o}所述，将第一个因子视为边界算子为 $\partial=\mu^1_{\A}$ 的链复形。由定义，
\begin{equation}\label{eq:5.15}
\hom_{\A}(\hom_{\A}(Y_0,Y_1)\otimes Y_0,Y_1)=\hom_{\A}(Y_0,Y_1)^{\vee}\otimes\hom_{\A}(Y_0,Y_1).
\end{equation}
这个态射群上的 $\mu^1_{\A}$ 算子为 $(c^{\vee}\otimes b)\mapsto(-1)^{|b|}\partial^{\vee}(c^{\vee})\otimes b+c^{\vee}\otimes\partial(b)$，其中 $\partial^{\vee}$ 为通常意义下的对偶微分。特别地，若把\eref{5.15}与 $\End_{\K}(\hom_{\A}(Y_0,Y_1))$ 等同，使 $c^{\vee}\otimes b$ 对应于自同态 $a\mapsto\langle c^{\vee},a\rangle b$，那么恒等自同态就是一个上闭链。将\eref{5.15}中相应的元素记为 $f$，并取\secref{3p}定义的映射锥 $\Cone(f)$。在 Yoneda 嵌入下，$\hom_{\A}(Y_0,Y_1)\otimes Y_0$ 成为 $\hom_{\A}(Y_0,Y_1)\otimes\Y_0$；$\mathcal I^1(f)$ 是\eref{5.2}定义的抽象求值映射 $\phi$；因此 $\mathcal I(\Cone(f))=\mathcal T_{Y_0}(\Y_1)$。\znote{71}{这里把锥 $\Cone(f)$ 显式写成扭曲复形，是扭曲复形一般定义的直接应用：对象为 $\hom_{\A}(Y_0,Y_1)\otimes Y_0$ 与 $Y_1$ 的直和，前者按基 $b^i$ 分解成若干 $Y_0$ 的平移，微分为 (5.18)，其下三角形状来自 $Y_0$ 与 $Y_1$ 的次序。一般框架中两对象扭曲复形的条件 $dq_{ij}+\sum_k q_{kj}q_{ik}=0$ 在这里恰好就是 $f$ 与 $\mu^1_{\A}$ 的相容性。[来源：nLab "twisted complex"，https:/\allowbreak /\allowbreak ncatlab.org/\allowbreak nlab/\allowbreak show/\allowbreak twisted+\allowbreak complex（2026-11-15 访问）。]}这就是说，在这种 $\Ai$ 范畴 $\A$ 中，可以取
\begin{equation}\label{eq:5.16}
T_{Y_0}Y_1=\Cone(f).
\end{equation}
\begin{note}\label{st:5.9}
把这一点写得更显式一些可能很有帮助，至少在 $Y_0,Y_1$ 是 $\B$ 的对象这一简单情形下如此。选取 $\hom_{\B}(Y_0,Y_1)$ 的一组基 $\{b^i\}$，指标为 $i\in I$，并写作 $\mu^1_{\B}(b^i)=\sum_j\beta^{ji}b^j$。那么 $T_{Y_0}(Y_1)$ 是扭曲复形
\begin{equation}\label{eq:5.17}
Y=\bigoplus_{i\in I}S^{-|b^i|+1}Y_0\oplus Y_1,
\end{equation}
其微分 $\delta_Y=(\phi^{ji}\otimes a^{ji})$ 中，$\phi$ 是 $\K$ 的各个平移之间的恒等映射，而
\begin{equation}\label{eq:5.18}
a^{ji}=\begin{cases}
(-1)^{\sigma^i+1}\beta^{ji}e_{Y_0},&i,j\in I,\\
(-1)^{\sigma^i}b^i,&i\in I,\ j=*,\\
0,&\text{其余情形}.
\end{cases}
\end{equation}
这里 $\sigma^i=|b^i|-1$，$*$ 是为\eref{5.17}中的 $Y_1$ 直和项添加的行、列所用的指标（或者，按照\stref{笔记}{3.26}修改等同\eref{3.29}后，\eref{5.18}两行中的符号 $(-1)^{\sigma^i}$ 都消失，分别剩下 $-\beta^{ji}e_{Y_0}$ 与 $b^i$）。
\end{note}
典范代表\eref{5.16}的存在具有许多实际优点。例如，可以直接写出公式来定义 $\Ai$ 函子 $T_{Y_0}$，而无须采用\secref{5b}中的间接构造。
\SourcePage{80}
\begin{note}\label{st:5.10}
设 $\widetilde{\A}$ 是 c-有限、三角化的 $\Ai$ 范畴。应用扰动引理，将其替换为 $\mu^1=0$ 的拟同构范畴，后者在链层次的 $\hom$ 空间因此有限维。再利用\stref{引理}{2.1}使之严格含幺。将结果记为 $\B$，并令 $\A=\Tw\B$。由构造，$\B$ 与 $\widetilde{\A}$ 拟同构，因此也是三角化的 $\Ai$ 范畴；这又意味着 $\A$ 与 $\widetilde{\A}$ 拟等价。利用这一点，常常可以把 c-有限、三角化 $\Ai$ 范畴的一般情形，归结为刚才讨论的技术上更简单的情形。
\end{note}

\subsection{伴随扭转}\label{sec:5g}
扭转运算有一个对偶，更准确地说是伴随运算。原则上它同样重要，但本书很少用到，因此这里只给出定义。取 $\Ai$ 范畴 $\A$ 的对象 $Y_0,Y_1$，如下定义 $\Ai$ 模 $\C^{\vee}$：
\begin{definition*}
\begin{equation}\label{eq:5.19}
\begin{aligned}
\C^{\vee}(X)&=\hom_{\A}(X,Y_1)\oplus\Hom_{\K}(\hom_{\A}(Y_1,Y_0),\hom_{\A}(X,Y_0))[-1],\\
\mu^1_{\C^{\vee}}(c,\beta)&=\bigl(\mu^1_{\A}(c),\ \mu^1_{\A}\circ\beta+(-1)^{|\beta|-1}\beta\circ\mu^1_{\A}+\mu^2_{\A}(\cdot,c)\bigr),\\
\mu^d_{\C^{\vee}}((c,\beta),a_{d-1},\ldots,a_1)
&=\bigl(\mu^d_{\A}(c,a_{d-1},\ldots,a_1),\\[-2pt]
&\qquad\mu^d_{\A}(\beta(\cdot),a_{d-1},\ldots,a_1)+\mu^{d+1}_{\A}(\cdot,c,a_{d-1},\ldots,a_1)\bigr).
\end{aligned}
\end{equation}
任意一个拟表示 $\C^{\vee}$ 的对象，都称为\emph{沿 $Y_0$ 对 $Y_1$ 作伴随扭转}的结果，记为 $T_{Y_0}^{\vee}Y_1$。
\end{definition*}
现在假设 $\A$ c-有限且三角化。那么，对所有 $Y_0,Y_1$，这样的拟表示对象都存在。它位于正合三角
\SourceEquationDiagram{diagram-5-20}{0.43615338}{5.20}
之中，其中的“对偶求值映射” $[f^{\vee}]$ 同样是所有可能的态射 $Y_1\to Y_0$ 的组合。

\subsection{球面对象}\label{sec:5h}
设 $\A$ 是 c-有限、三角化的 $\Ai$ 范畴。沿用最初由 Kontsevich 提出的想法，作如下定义：
\begin{definition*}
若 $Y_0\in\Ob\A$ 满足以下条件，就称它为维数 $n\in\Z$ 的\emph{球面对象}。
\begin{itemize}
\item （Poincaré 对偶性）存在一个线性映射 $\int:\Hom^n_{H(\A)}(Y_0,Y_0)\to\K$，使对所有 $Y_1$，配对
\[
\Hom^{n-\bullet}_{H(\A)}(Y_1,Y_0)\otimes\Hom^{\bullet}_{H(\A)}(Y_0,Y_1)
\xrightarrow{\text{复合}}\Hom^n_{H(\A)}(Y_0,Y_0)\xrightarrow{\int}\K
\]
非退化。
\SourcePage{81}
\item （球面自同态环）$Y_0$ 的自同态空间满足下列短正合列
\[
0\longrightarrow\K\xrightarrow{[e_{Y_0}]}\Hom_{H(\A)}(Y_0,Y_0)\xrightarrow{\int}\K\longrightarrow0.
\]
\end{itemize}
\end{definition*}
若 $n\ne0$，最后一个条件等价于 $\Hom_{H(\A)}(Y_0,Y_0)\cong\K[t]/t^2$，其中 $t$ 的次数为 $n$。当 $n>0$ 时，这恰好是 $S^n$ 的通常上同调环，这也解释了所用术语。后面将看到，Lagrangian 球面（以及同调球面）在 Fukaya 范畴中成为球面对象。\znote{73}{“球面对象”的术语来自拓扑：当维数 $n>0$ 时，自同态环 $\K[t]/t^2$（$\deg t=n$）与球面 $S^n$ 的上同调环 $\HH^*(S^n)$ 一致，而 Poincaré 对偶条件则模仿了上同调中的相交配对。在 Fukaya 范畴中，精确 Lagrangian 球面 $\mathbb{S}^n\subset\mathbb{C}^{n+1}$ 的 Floer 上同调同构于 $\HH^*(S^n)$，配对由 Lagrangian 相交数给出，因此它确实成为球面对象。[来源：译者按（上下文推断，无外部依据）：由本节定义与 $\HH^*(S^n)$ 的经典计算推得。]}
\begin{lemma}\label{st:5.11}
若 $Y_0$ 是球面对象，则 $T_{Y_0}$ 是从 $\A$ 到自身的拟等价。
\end{lemma}
\begin{proof}
给定 $\A$ 的两个对象 $Y_k$（$k=1,2$），令 $C$ 为映射 $\mathcal T_{Y_0}^1:\hom_{\Q}(\Y_1,\Y_2)\to\hom_{\Q}(\mathcal T_{Y_0}\Y_1,\mathcal T_{Y_0}\Y_2)$ 的映射锥（按通常链复形的意义）。它带有三步递降滤过 $F^{\bullet}C$，其相伴分次空间为
\[
\begin{aligned}
C/F^1C&=\hom_{\Q}(\Y_1,\hom_{\A}(Y_0,Y_2)\otimes\Y_0)[1],\\
F^1C/F^2C&=\hom_{\Q}(\Y_1,\Y_2)[1]\oplus\hom_{\Q}(\Y_1,\Y_2)\\[-2pt]
&\qquad\oplus\hom_{\Q}(\hom_{\A}(Y_0,Y_1)\otimes\Y_0,\hom_{\A}(Y_0,Y_2)\otimes\Y_0),\\
F^2C&=\hom_{\Q}(\hom_{\A}(Y_0,Y_1)\otimes\Y_0,\Y_2)[-1].
\end{aligned}
\]
由于 $F^1C/F^2C$ 层次上的抵消，所得谱序列具有以下非零项：
\[
E_1^{rs}=\begin{cases}
(\Hom_{H(\A)}(Y_0,Y_2)\otimes\Hom_{H(\A)}(Y_1,Y_0))^{r+s+1},&r=0,\\
\bigl(\Hom_{H(\A)}(Y_0,Y_2)\otimes\Hom_{H(\A)}(Y_0,Y_0)\\
\qquad\otimes\Hom_{H(\A)}(Y_0,Y_1)^{\vee}\bigr)^{r+s},&r=1,\\
(\Hom_{H(\A)}(Y_0,Y_2)\otimes\Hom_{H(\A)}(Y_0,Y_1)^{\vee})^{r+s-1},&r=2.
\end{cases}
\]
$E_1$ 页的微分 $\partial^{rs}$ 由 $[\phi_k]$ 诱导，也就是说，它由 $H(\A)$ 中的复合映射或其对偶组成。例如，$\partial^{1\bullet}$ 满射，因为
\SourceDiagram{diagram-5-11a}{0.55661479}
\SourcePage{82}
是恒等映射。类似地，$\partial^{0\bullet}$ 单射，因为与 Poincaré 映射复合后得到同构：
\SourceDiagram{diagram-5-11b}{0.55661479}
再作维数计算，就可见谱序列的 $E_2$ 页为零。由构造，这意味着 $T_{Y_0}$ 上同调全忠实。\znote{74}{证明的策略是：把 $T_{Y_0}$ 在 $\hom$ 空间上的作用看成对映射锥 $C$ 作三步滤过，$E_1$ 页由三项组成，分别涉及 $\Hom(Y_0,Y_2)$、$\Hom(Y_0,Y_2)\otimes\Hom(Y_0,Y_0)\otimes\Hom(Y_0,Y_1)^{\vee}$ 与 $\Hom(Y_0,Y_2)\otimes\Hom(Y_0,Y_1)^{\vee}$。Poincaré 对偶使 $E_1$ 页的两个微分一满一单，配合有限维性作维数计数即得 $E_2=0$；等价地，$T_{Y_0}$ 在同调层次上全忠实。[来源：并见 Seidel, Fukaya categories and Picard–Lefschetz theory, §5h（引理 5.11 的证明）。]}

还需证明它在对象的同构类上满射。用 $\B\subset\A$ 表示由同构于某个 $T_{Y_0}(Y_1)$ 的对象组成的全 $\Ai$ 子范畴。由\eref{5.4}与 $Y_0$ 的球面性可得
\[
T_{Y_0}(S^{n-1}Y_0)\cong Y_0,
\]
因此 $Y_0$ 属于 $\B$。利用这一点，再用\eref{5.4}，可见 $\B$ 生成 $\A$。另一方面，\stref{推论}{3.9}表明，$\B$ 本身是三角化的 $\Ai$ 范畴，故包含函子 $\B\hookrightarrow\A$ 必为拟等价。
\end{proof}

\begin{note}\label{st:5.12}
事实上，在 $H^0(\operatorname{fun}(\A,\A))$ 中的同构意义下，伴随扭转 $T_{Y_0}^{\vee}$ 是 $T_{Y_0}$ 的逆。这里不证明这一点（经典导出范畴背景下的平行定理，见\bcite{134}命题 2.10），但要说明一个部分结果，即在 $H^0(\A)$ 中有
\begin{equation}\label{eq:5.21}
T_{Y_0}T_{Y_0}^{\vee}(Y_1)\cong Y_1.
\end{equation}
为此，转到模范畴 $\Q=\modcat(\A)$，取\eref{5.19}中的 $\C^{\vee}$。$\C^{\vee}(Y_0)$ 的上同调位于长正合列
\begin{equation}\label{eq:5.22}
\begin{gathered}
\vdots\\
\Hom_{H(\A)}(Y_0,Y_1)[-1]\\
\big\downarrow\\
\Hom_{H(\A)}(Y_1,Y_0)^{\vee}\otimes\Hom_{H(\A)}(Y_0,Y_0)[-1]\\
\big\downarrow\\
H(\C^{\vee}(Y_0))\\
\vdots
\end{gathered}
\end{equation}
之中，其中第一个同态又是一个对偶复合映射。由 Poincaré 对偶性，经过投影 $\int:\Hom_{H(\A)}(Y_0,Y_0)\to\K$ 后，它成为同构。利用
\SourcePage{83}
$Y_0$ 的自同态环的球面性，可知\eref{5.22}中第二个映射在子空间 $\Hom_{H(\A)}(Y_1,Y_0)^{\vee}\otimes\K\cdot[e_{Y_0}]$ 上的限制为同构。因此，将 $\mathcal T_{Y_0}\C^{\vee}$ 构造为抽象求值映射的锥时，可以把 $\C^{\vee}(Y_0)$ 替换为拟同构的子复形 $\hom_{\A}(Y_1,Y_0)^{\vee}=\Hom_{\K}(\hom_{\A}(Y_1,Y_0),\K\cdot e_{Y_0})$。替换后得到的 $\Ai$ 模为
\[
\begin{aligned}
\M(X)&=(\hom_{\A}(Y_1,Y_0)^{\vee}\otimes\hom_{\A}(X,Y_0))\oplus\hom_{\A}(X,Y_1)\\[-2pt]
&\qquad\oplus\Hom_{\K}(\hom_{\A}(Y_1,Y_0),\hom_{\A}(X,Y_0))[-1],\\
\mu^1_{\M}(\gamma\otimes b,c,\beta)
&=\bigl((-1)^{|b|+|\gamma|}(\gamma\circ\mu^1_{\A})\otimes b+\gamma\otimes\mu^1_{\A}(b),\ \mu^1_{\A}(c),\\[-2pt]
&\qquad\mu^1_{\A}\circ\beta+(-1)^{|\beta|-1}\beta\circ\mu^1_{\A}\\[-2pt]
&\qquad+\mu^2_{\A}(\cdot,c)+\mu^2_{\A}(e_{Y_0},b)\otimes\gamma(\cdot)\bigr),\\
\mu^d_{\M}((\gamma\otimes b,c,\beta),a_{d-1},\ldots,a_1)
&=\bigl(\gamma\otimes\mu^d_{\A}(b,a_{d-1},\ldots,a_1),\ \mu^d_{\A}(c,a_{d-1},\ldots,a_1),\\[-2pt]
&\qquad\mu^d_{\A}(\beta(\cdot),a_{d-1},\ldots,a_1)+\mu^{d+1}_{\A}(\cdot,c,a_{d-1},\ldots,a_1)\\[-2pt]
&\qquad+\mu^{d+1}_{\A}(e_{Y_0},b,a_{d-1},\ldots,a_1)\otimes\gamma(\cdot)\bigr).
\end{aligned}
\]
存在模同态 $\M\to\Y_1$，其一阶项为 $(\gamma\otimes b,c,\beta)\mapsto(-1)^{|c|}c$，高阶项为零。它的核是 $\M(X)$ 中满足 $c=0$ 的子复形；利用 $e_{Y_0}$ 在上同调层次上是恒等态射这一事实，容易看出该核无上同调。因此，在 $H^0(\Q)$ 中有同构 $\mathcal T_{Y_0}\C^{\vee}\cong\M\cong\Y_1$。利用定义回到 $H^0(\A)$，就得到\eref{5.21}。
\end{note}
